\chapter{THE EXPONENTIAL FUNCTION}

The most widely known result on transcendental numbers is the transcendency of $e$ proved by Lindemann in 1882. His method is based on Hermite's previous work who discovered the transcendency of $e$ in 1873. Both results are contained in the general Lindemann-Weierstrass theorem which will be proved in \S 12. We shall start with some simpler problems, namely the irrationality of $e$ and related questions.

\section{The irrationality of $e$}

The usual proof of the irrationality of $e$ runs as follows. From the series
\[ e = \sum_{k=0}^{\infty} \frac{1}{k!} \]
we get the decomposition
\[ e = s_n + r_n, \quad s_n = \sum_{k=0}^{n} \frac{1}{k!}, \quad r_n = \sum_{k=n+1}^{\infty} \frac{1}{k!} \quad (n=1, 2, \dots). \]
Since
\[ r_n = \frac{1}{(n+1)!} \left( 1 + \frac{1}{n+2} + \frac{1}{(n+2)(n+3)} + \dots \right) < \frac{e-1}{(n+1)!} \]
we find that
\[ e = s_1 + r_1 < 2 + \frac{e-1}{2}, \quad e < 3; \]
therefore
\[ 0 < r_n < \frac{2}{(n+1)!}. \]
Put
\[ n! s_n = a_n, \quad n! r_n = b_n, \]
then the number $a_n$ is integral and
\[ 0 < b_n < \frac{2}{n+1} \le 1 \]
for $n=1, 2, \dots$. This proves that $n! e = a_n + b_n$ and, a fortiori, $n! e$ is never an integer. In other words, $e$ is irrational.

The proof is still simpler if we use the series for $e^{-1}$ instead of $e$. Then
\[ e^{-1} = \sigma_n + \rho_n, \quad \sigma_n = \sum_{k=0}^{n} \frac{(-1)^k}{k!}, \quad \rho_n = \sum_{k=n+1}^{\infty} \frac{(-1)^k}{k!} \quad (n=1, 2, \dots) \]
and
\[ 0 < (-1)^{n+1} \rho_n = \frac{1}{(n+1)!} - \frac{1}{(n+2)!} + \dots < \frac{1}{(n+1)!}. \]
Defining
\[ n! \sigma_n = \alpha_n, \quad n! \rho_n = \beta_n, \]
we see that $\alpha_n$ is integral and
\[ 0 < (-1)^{n+1} \beta_n < \frac{1}{n+1} < 1. \]
Therefore $n! e^{-1} = \alpha_n + \beta_n$ and, a fortiori, $n! e^{-1}$ is never an integer.

We can prove a little more, namely that $e$ is not the root of a quadratic equation $ax^2 + bx + c = 0$ with integral $a, b, c$, not all 0. Consider the expression
\[ E_n = n! (ae + ce^{-1}) \]
with integral $a$ and $c$, not both 0. Then
\[ E_n = S_n + R_n, \quad S_n = a a_n + c \alpha_n, \quad R_n = a b_n + c \beta_n \]
where $S_n$ is integral and the absolute value
\[ |R_n| \leq |ab_n| + |c\beta_n| < \frac{2|a|+|c|}{n+1}, \]
so that
\[ |R_n| < 1 \]
for all $n \ge 2|a|+|c|$. On the other hand we have the recursion formula
\[ \begin{aligned} R_{n-1} - R_n &= a(nb_{n-1} - b_n) + c(n\beta_{n-1} - \beta_n) \\ &= a + (-1)^{n}c. \end{aligned} \]
It follows that at least one of the three numbers $R_{n-1}$, $R_n$, $R_{n+1}$ is different from 0, since otherwise $a+c=0, a-c=0$ and $a=0, c=0$. This shows the existence of a positive integer $\nu$ such that $E_{\nu}$ is not integral, and therefore the number
\[ b + \frac{E_{\nu}}{\nu!} = ae + b + ce^{-1} \]
is different from 0, for all integral $b$. This means that
\[ ae^2 + be + c \neq 0 \]
for arbitrary integers $a, b, c$, not all 0. In other words, $e$ is not a quadratic irrationality.

\section{The operator $f(D)$}

We denote by $D$ the differentiation with respect to the variable $x$. If
\[ f(t) = a_0 + a_1 t + a_2 t^2 + \dots \]
is a power series with real or complex coefficients $a_0, a_1, \dots$ and $\varphi = \varphi(x)$ a function of $x$, we define
\begin{equation} f(D)\varphi = \sum_{n=0}^{\infty} a_n D^n \varphi = \sum_{n=0}^{\infty} a_n \frac{d^n \varphi}{dx^n}. \tag{1}\label{eq:1} \end{equation}
In order to avoid questions of convergence and differentiability, we shall apply the operator $f(D)$ only in two cases: either $\varphi$ is a polynomial, or $f$ is a polynomial and $\varphi$ possesses derivatives of all orders. In both cases, the series\eqref{eq:1} is finite.

It is clear that, for two power series $f_1(t)$ and $f_2(t)$,
\begin{equation} f_1(D)(f_2(D)\varphi) = (f_1 f_2)(D)\varphi \tag{2}\label{eq:2} \end{equation}
Moreover, if $a_0 \neq 0$, then
\[ f^{-1}(t) = b_0 + b_1 t + \dots \quad (b_0 = a_0^{-1}) \]
exists and, by\eqref{eq:2},
\[ f^{-1}(D)(f(D)\varphi) = \varphi \]
for polynomials $\varphi$.

If the power series $f$ and the polynomial $\varphi$ have integral coefficients, then also $f(D)\varphi$ does, and the same holds for $f^{-1}(D)\varphi$, in case $a_0 = \pm 1$.

We define
\[ J \varphi = \int_0^x \varphi(t) dt \]
so that
\[ D J \varphi = \varphi(x), \qquad J D \varphi = \varphi(x) - \varphi(0) \]
and
\begin{equation} J^{n+1} \varphi = \int_0^x \frac{(x-t)^n}{n!} \varphi(t) dt \quad (n = 0,1,\dots). \tag{3}\label{eq:3} \end{equation}
In particular we are interested in the case $\varphi = e^{\lambda x} P$, where $\lambda$ is a constant and $P = P(x)$ a polynomial. Since
\[ D(e^{\lambda x} P) = e^{\lambda x} (\lambda + D)P, \]
we have
\begin{equation} D^{n}(e^{\lambda x}P) = e^{\lambda x}(\lambda + D)^{n}P \qquad (n=1,2,\dots) \tag{4}\label{eq:4} \end{equation}
and
\begin{equation} (\lambda + D)^{n}P = Q(x) \tag{5}\label{eq:5} \end{equation}
again is a polynomial. Vice versa, if $\lambda \neq 0$ and a polynomial $Q$ are arbitrarily given, then\eqref{eq:5} implies
\[ P = (\lambda + D)^{-n}Q \]
and this is the unique polynomial solution of the differential equation
\[ D^{n}(e^{\lambda x}P) = e^{\lambda x}Q. \]

In case $\lambda = \pm 1$, the polynomial $P$ has integral coefficients if $Q$ has.

\section{Approximation to $e^x$ by rational functions}

We are going to determine two polynomials $A = A(x)$, $B = B(x)$ of degree $n$, such that the sum $B e^x + A$ vanishes at $x = 0$ of order $2n + 1$. This condition implies
\begin{equation} B e^x + A = R = c x^{2n+1} + \dots \tag{6}\label{eq:6} \end{equation}
where $R = R(x)$ is a power series starting with a term of order $2n + 1$. Writing $A$ and $B$ with indeterminate coefficients and equating the terms of order $0, 1, \dots, 2n$ in\eqref{eq:6}, we obtain $2n+1$ homogeneous linear equations for the $2n+2$ unknown coefficients in $A$ and $B$. This proves that\eqref{eq:6} has a non-trivial solution $A, B$. It will turn out that $c \neq 0$.

In order to obtain an explicit formula for $A$ and $B$, we differentiate\eqref{eq:6} $n+1$ times; then
\[ D^{n+1}(B e^x) = D^{n+1}R \]
\begin{equation} e^x(1+D)^{n+1}B = D^{n+1}R = c_0 x^n + \dots \tag{7}\label{eq:7} \end{equation}
where
\[ c_0 = (2n+1) \dots (n+1)c. \]
Therefore
\begin{equation} (1+D)^{n+1}B = e^{-x}(c_0 x^n + \dots) = c_0 x^n + \dots = c_0 x^n \tag{8}\label{eq:8} \end{equation}
because $(1+D)^{n+1}B$ is a polynomial of degree $n$ at most, and
\[ B = c_0(1+D)^{-n-1}x^n. \]

To find $A$, we get in the same manner
\[ D^{n+1}(A e^{-x}) = D^{n+1}(R e^{-x}) \]
\[ e^{-x}(-1+D)^{n+1}A = c_0 x^n + \dots \]
\[ (-1+D)^{n+1}A = e^x(c_0 x^n + \dots) = c_0 x^n + \dots = c_0 x^n \]
\[ A = c_0(-1+D)^{-n-1}x^n. \]
This proves that $A$ and $B$ are unique, up to the arbitrary constant factor $c_0$. Choose $c_0 = 1$, then
\begin{equation} A(x) = (-1+D)^{-n-1}x^n, \quad B(x) = (1+D)^{-n-1}x^n \tag{9}\label{eq:9} \end{equation}
have integral coefficients. Moreover, by\eqref{eq:7} and\eqref{eq:8},
\[ D^{n+1}R = x^n e^x. \]
Since $R$ and its first $n$ derivatives are $0$ at $x = 0$, we obtain, by\eqref{eq:3},
\[ R = J^{n+1}D^{n+1}R = \frac{1}{n!} \int_{0}^{x} (x-t)^{n}t^{n}e^{t}dt \]
\begin{equation} R(x) = \frac{x^{2n+1}}{n!} \int_0^1 t^n (1-t)^n e^{tx} dt \tag{10}\label{eq:10} \end{equation}
for $n=0,1,\dots$. Replacing $t$ by $1-t$ we find
\[ R(x) = \frac{x^{2n+1}}{n!} \int_0^1 t^n (1-t)^n e^{(1-t)x} dt \]
so that
\begin{equation} R(x) = \frac{x^{2n+1}}{n!} e^{\frac{x}{2}} \int_{0}^{1} t^{n} (1-t)^{n} \cosh\left\{\left(t-\frac{1}{2}\right)x\right\} dt \quad (n = 0, 1, \dots). \tag{11}\label{eq:11} \end{equation}

\section{The irrationality of $e^a$ for rational $a \neq 0$}

By\eqref{eq:10},
\begin{equation} |R(x)| \leq \frac{|x|^{2n+1}}{n!} e^{|x|}, \tag{12}\label{eq:12} \end{equation}
for all complex $x$, and
\begin{equation} R(x) > 0, \tag{13}\label{eq:13} \end{equation}
for all positive $x$.

Now let $x = m$ be a positive integer; then the numbers $A(m)$ and $B(m)$ are integers. Suppose that $e^m$ were rational and denote by $q > 0$ the denominator of $e^m$. By\eqref{eq:6}, the number
\[ q R(m) = r \]
is integral. Because of\eqref{eq:12} and\eqref{eq:13},
\[ 0 < r \leq q \frac{m^{2n+1}}{n!} e^m = q m e^m \frac{m^{2n}}{n!} < 1 \]
for all sufficiently large $n$. This is a contradiction.

Therefore all powers $e^m$ ($m=1,2,\dots$) are irrational. If $a$ is any rational number $\neq 0$, we write $a=\frac{m}{r}$ with integral $m>0, r \neq 0$. Since $e^m=(e^a)^r$, it follows that $e^a$ is irrational, for all rational $a \neq 0$.

\section{The irrationality of $\pi$}

We know that the polynomials $A(x)$ and $B(x)$ of degree $n$ are uniquely determined by the formula
\begin{equation} B(x)e^{x} + A(x) = R(x) = cx^{2n+1} + \dots, \tag{14}\label{eq:14} \end{equation}
where $c \neq 0$ is given. Replace $x$ by $-x$ and multiply by $e^{x}$; then
\[ A(-x)e^{x} + B(-x) = e^{x}R(-x) = -cx^{2n+1} + \dots, \]
whence
\begin{equation} A(-x) = -B(x). \tag{15}\label{eq:15} \end{equation}
This can also be proved by using the expressions\eqref{eq:9} for $A$ and $B$.

Choose $x=\pi i$ and apply\eqref{eq:11}, \eqref{eq:14} and\eqref{eq:15}; then
\[ e^{\frac{x}{2}} = i, \quad \cosh \left\{\left(t-\frac{1}{2}\right)x\right\} = \cos \left\{\left(t-\frac{1}{2}\right)\pi\right\} = \sin \pi t \]
\begin{equation} A(\pi i) + A(-\pi i) = R(\pi i) \tag{16}\label{eq:16} \end{equation}
and
\begin{equation} R(\pi i) = (-1)^{n+1} \frac{\pi^{2n+1}}{n!} \int_0^1 t^n (1-t)^n \sin \pi t \, dt. \tag{17}\label{eq:17} \end{equation}
The integrand is positive in the interval $0 < t < 1$, so that
\[ R(\pi i) \neq 0. \]
The function $A(x) + A(-x)$ is a polynomial of degree $\nu = \left[\frac{n}{2}\right]$ in the variable $x^2$ with integral coefficients. If $\pi^2$ is a rational number and $q > 0$ its denominator, then the number
\begin{equation} q^{\nu} R(\pi i) = j \tag{18}\label{eq:18} \end{equation}
is integral, by\eqref{eq:16}. However, it follows from\eqref{eq:17} and\eqref{eq:18} that
\[ 0 < |j| \leq \frac{q^{\frac{n}{2}} \pi^{2n+1}}{n!} < 1 \]
for all sufficiently large values of $n$.

This contradiction proves that $\pi^2$ and, a fortiori, $\pi$ itself are irrational.

\section{The irrationality of $\tan a$ for rational $a \ne 0$}

Define
\[ A(x) - A(-x) = x P(x^2), \quad A(x) + A(-x) = Q(x^2), \]
so that $P=P(x^2), Q=Q(x^2)$ are polynomials in $x^2$ of degrees $[\frac{n-1}{2}]$, $[\frac{n}{2}] = \nu$, with integral coefficients, and
\[ 2A(x) = Q+xP, \quad 2A(-x) = Q-xP. \]
Multiply\eqref{eq:14} by $e^{-\frac{x}{2}}$ and use\eqref{eq:11}, \eqref{eq:15}; then
\begin{equation} \begin{aligned} A(x) e^{-\frac{x}{2}} - A(-x)e^{\frac{x}{2}} &= R(x) e^{-\frac{x}{2}} \\ xP \cosh \frac{x}{2} - Q \sinh \frac{x}{2} &= R(x)e^{-\frac{x}{2}} \\ &= \frac{x^{2n+1}}{n!} \int_{0}^{1} t^{n} (1-t)^{n} \cosh\left\{\left(t-\frac{1}{2}\right)x\right\} dt \\ &= S(x), \end{aligned} \tag{19}\label{eq:19} \end{equation}
say.

Now let $a^2=\pm b$ be a rational number $\neq 0$ and assume that also $\frac{\tanh a}{a}$ is rational. Put $x=2a$ and denote by $q > 0$ the denominator of $x^2=4a^2=\pm 4b$. Then
\[ x\cosh \frac{x}{2}=\gamma r, \quad \sinh \frac{x}{2}=\gamma s \]
where $r, s$ are integers and $\gamma \neq 0$, the numbers $q^\nu P(x^2), q^\nu Q(x^2)$ are integral, and\eqref{eq:19} shows that also 
\[ \gamma^{-1} q^{\nu} S(x) = j \]
is an integer, whereas
\[ |j| \leq \left| \frac{x}{\gamma} \right| e^{\frac{|x|}{2}}  \frac{q^{\frac{n}{2}}|x|^{2n}}{n!} \longrightarrow 0 \quad (n \longrightarrow \infty). \]
We obtain the desired contradiction if we can prove that $R(x) \neq 0$. This inequality is obvious in case $x^2 \geq -\pi^2$, since then the integrand in\eqref{eq:19} is positive for $0 < t < 1$. In order to complete the proof in the remaining case, we write more explicitly $A=A_n$, $B=B_n$, $R=R_n$, $c=c_n$; then
\[ \begin{aligned} B_n e^{x} + A_n &= R_n = c_n x^{2n+1} + \dots, \\ B_{n-1} e^{x} + A_{n-1} &= R_{n-1} = c_{n-1} x^{2n-1} + \dots, \quad (n = 1, 2, \dots) \end{aligned} \]
\begin{equation} \begin{aligned} A_{n-1} B_n - A_n B_{n-1} &= R_{n-1} B_n - R_n B_{n-1} \\ &= c_{n-1} B_n(0) x^{2n-1} + \dots \\ &= c_{n-1} B_n(0) x^{2n-1}, \end{aligned} \tag{20}\label{eq:20} \end{equation}
because $A_{n-1} B_n - A_n B_{n-1}$ is a polynomial in $x$ of degree $2n-1$. If $B_n(0) = 0$, then also $A_n(0) = B_n(0) - R_n(0) = 0$, and the formula
\[ (B_{n-1} + x^{-1}B_n)e^{x} + (A_{n-1} + x^{-1}A_n) = c_{n-1}x^{2n-1} + \dots \]
would contradict the unicity of $A_{n-1}$, $B_{n-1}$. Therefore $B_n(0) \neq 0$. This also follows from $B(x) = (1+D)^{-n-1}x^n$, namely $B_{n}(0)= \binom{-n-1}{n}n! \neq 0$.

Now we infer from\eqref{eq:20} that $R_{n-1}B_n - R_nB_{n-1} \neq 0$ for all $x \neq 0$, so that at least one of the two numbers $R_{n-1}(x)$, $R_n(x)$ is different from 0. This suffices for the completion of the proof.

Separating the real and imaginary case, we see that $\frac{\tanh \sqrt{b}}{\sqrt{b}}$ and $\frac{\tan \sqrt{b}}{\sqrt{b}}$ are irrational, for all rational positive numbers $b$. In particular, $\tan a$ is irrational for all rational $a \neq 0$. This contains the irrationality of $\pi$, since $\tan \frac{\pi}{4} = 1$ is rational.

\begin{remark} note：Siegel use $\text{tg}(x)$ to express $\tan(x)$, we use the morden notation. \end{remark}

\section{The function $P_1 e^{\rho_1 X} + \dots + P_m e^{\rho_m X}$}

We are going to solve the following problem: Let $m$ different complex constants $\rho_1,\ldots,\rho_m$ and $m$ nonnegative integers $n_1,\ldots,n_m$ be given and put
\begin{equation} \sum_{k=1}^{m} (n_k + 1) = N + 1 \tag{21}\label{eq:21} \end{equation}
to determine $m$ polynomials $P_1(x),\ldots, P_m(x)$ of degrees $n_1,\ldots,n_m$ such that the function
\begin{equation} R = P_{1}e^{\rho_{1}x} + \dots + P_{m}e^{\rho_{m}x} \tag{22}\label{eq:22} \end{equation} 
vanishes at $x=0$ of order $N$. For the special case $m=2$, $n_1=n_2=n$, $\rho_1=1$, $\rho_2=0$, the solution of this problem was given in \S 3.

Writing $P_1, \dots, P_m$ with indeterminate coefficients and considering the terms of degree $0, 1, \dots, N-1$ in the power series expansion of $R$, we obtain $N$ homogeneous linear equations for the $(n_1+1) + \dots + (n_m+1)$ unknown coefficients. These equations have a non-trivial solution, because of\eqref{eq:21}. This proves the existence of $m$ polynomials $P_k$ ($k=1,\dots,m$), not all identically 0, of degrees $\leq n_k$, such that the power series of $R$ takes the form
\begin{equation} R = c \frac{x^{N}}{N!} + \dots \tag{23}\label{eq:23} \end{equation}
with some constant $c$.

In case $m=1$ the solution clearly is
\[ P_1 = c \frac{x^{n_1}}{n_1!} , \]
so that $c \neq 0$. It will turn out that, for every $m$, the constant $c$ in\eqref{eq:23} is different from 0 and that, for any given $c \neq 0$, the solution is unique.

Write more explicitly $N=N_m$ and suppose $m > 1$. By\eqref{eq:4} and\eqref{eq:22},
\begin{equation} \begin{aligned} D^{n_m+1}(Re^{-\rho_m x}) &= \sum_{k=1}^{m} D^{n_m+1}(e^{(\rho_k-\rho_m)x}P_k) \\ &= \sum_{k=1}^{m-1} e^{(\rho_k-\rho_m)x}(\rho_k-\rho_m+D)^{n_m+1}P_k; \end{aligned} \tag{24}\label{eq:24} \end{equation}
by\eqref{eq:23},
\begin{equation} \begin{aligned} D^{n_m+1}(Re^{-\rho_m x}) &= c D^{n_m+1}\left(\frac{x^{N_m}}{N_m!}e^{-\rho_m x}\right) \\ &= c\frac{x^{N_{m-1}}}{N_{m-1}!} + \dots \end{aligned} \tag{25}\label{eq:25} \end{equation}
Since the $m-1$ polynomials
\[ Q_k = (\rho_k - \rho_m + D)^{n_m+1} P_k \quad (k=1,\dots,m-1) \]
have exactly the same degrees as $P_k$, they do not all vanish identically. It follows from\eqref{eq:24} and\eqref{eq:25} that the function
\[ S = D^{n_m+1} (Re^{-\rho_m x}) = Q_1 e^{(\rho_1-\rho_m)x} + \dots + Q_{m-1} e^{(\rho_{m-1}-\rho_m)x} \]
solves the problem for $m-1$ and $\rho_1 - \rho_m, \ldots, \rho_{m-1} - \rho_m$ instead of $m$ and $\rho_1, \ldots, \rho_m$, with the same constant $c$. We now obtain
\[ P_k = (\rho_k - \rho_m + D)^{-n_m-1} Q_k \quad (k=1,\dots,m-1) \]
and, because of\eqref{eq:3},
\begin{equation} R = e^{\rho_m x} J^{n_m+1} S = e^{\rho_m x} \int_0^x \frac{(x-t)^{n_m}}{n_m!} S(t) dt. \tag{26}\label{eq:26} \end{equation}

Applying induction with respect to $m$, we see that we can prescribe $c=1$ and that then
\[ P_1 = \prod_{l=2}^{m} \left( \rho_1 - \rho_l + D \right)^{-n_l-1} \frac{x^{n_1}}{n_1!}, \]
more generally
\begin{equation} P_k = \prod_{\substack{l=1\\l \neq k}}^{m} (\rho_k - \rho_l + D)^{-n_l-1} \frac{x^{n_k}}{n_k!} \quad (k=1,\dots,m). \tag{27}\label{eq:27} \end{equation}

In order to determine $R$ explicitly we consider first the case $m=2$. By\eqref{eq:26},
\[ \begin{aligned} R &= e^{\rho_2 x} \int_0^x \frac{(x-t)^{n_2}}{n_2!} e^{(\rho_1 - \rho_2)t} \frac{t^{n_1}}{n_1!} dt \\ &= \int_{\substack{t_1 + t_2 = x \\ t_1 > 0, t_2 > 0}} \frac{t_1^{n_1} t_2^{n_2}}{n_1! n_2!} e^{\rho_1 t_1 + \rho_2 t_2} dt_1 \qquad (x>0). \end{aligned} \]
Suppose that the formula
\begin{equation} R(x) = \int \dots \int_{\substack{t_1 + \dots + t_m = x \\ t_1 > 0, \dots, t_m > 0}} \prod_{k=1}^{m} \left( \frac{t_k^{n_k}}{n_k!} e^{\rho_k t_k} \right) dt_1 \dots dt_{m-1} \quad (x>0) \tag{28}\label{eq:28} \end{equation}
is true for $m-1 \geq 1$, instead of $m$. Then
\[ S(t) = \int \dots \int_{\substack{t_1 + \dots + t_{m-1} = t \\ t_1 > 0, \dots, t_{m-1} > 0}} \prod_{k=1}^{m-1} \left( \frac{t_k^{n_k}}{n_k!} e^{(\rho_k - \rho_m)t_k} \right) dt_1 \dots dt_{m-2} \quad (t > 0). \]
Substituting in\eqref{eq:26} and defining $t_m=x-t$, we obtain\eqref{eq:28} for $m$.

We shall call $R = P_1 e^{\rho_1 x} + \dots + P_m e^{\rho_m x}$ an approximation form.

\section{Estimation of $R(1)$}

From\eqref{eq:28} we can draw two important conclusions: For arbitrary complex $\rho_1,\ldots,\rho_m$ we have the upper estimate
\begin{equation} |R(1)| \leq \frac{e^{|\rho_1| + \dots + |\rho_m|}}{n_1! \dots n_m!}, \tag{29}\label{eq:29} \end{equation}
and for \textit{real} $\rho_1, \dots, \rho_m$ we have the \textit{lower} estimate
\begin{equation} R(1) > 0. \tag{30}\label{eq:30} \end{equation}

\section{Estimation of $P_k(1)$ and its denominator}

We apply the expansion
\begin{equation} (\omega + D)^{-n-1} = \omega^{-n-1} \sum_{r=0}^{\infty} \binom{-n-1}{r} \omega^{-r} D^r \quad (\omega \neq 0) \tag{31}\label{eq:31} \end{equation}
to find an estimate of $P_k$ in\eqref{eq:27}. Denote by $M$ the maximum of the $m(m-1)$ numbers $\frac{1}{|\rho_k - \rho_l|}$ ($1 \leq k, l \leq m, k \neq l$). Then
\[ |P_{k}(x)| \le (M^{-1} - D)^{n_{k} - N} \frac{x^{n_{k}}}{n_{k}!} \quad (x>0), \]
since the right-hand side, as a polynomial in $x$, majorizes $P_k(x)$. Furthermore,
\[ \begin{aligned} (M^{-1} - D)^{n_k - N} \frac{x^{n_k}}{n_k!} &= M^{N - n_k} \sum_{r=0}^{n_k} \binom{N - n_k + r - 1}{r} M^r D^r \frac{x^{n_k}}{n_k!} \\ &\leq \sum_{r=0}^{n_k} \binom{N}{r} M^{N-n_k+r} \frac{x^{n_k-r}}{(n_k-r)!} \\ &\leq (1+1)^N (M+x)^N = (2M+2x)^N \quad (x>0), \end{aligned} \]
so that
\begin{equation} |P_{k}(1)| \le (2M+2)^{N} \quad (k=1,\dots,m). \tag{32}\label{eq:32} \end{equation}

If $\rho_1, \dots, \rho_m$ are algebraic numbers, then also $P_k(1)$ is algebraic. To find an estimate of the denominator of $P_k(1)$ we choose a positive rational integer $q$ such that the $m(m-1)$ algebraic numbers $q(\rho_k - \rho_l)^{-1}$ ($1 \leq k, l \leq m, k \neq l$) are all integral. It follows from\eqref{eq:27} and\eqref{eq:31}, that the number $n_k! q^N P_k(1)$ is integral.

\section{The transcendency of $e^a$ for real algebraic $a \neq 0$}

Let $a \neq 0$ be a real algebraic number and suppose that also $e^a$ were algebraic. We introduce the algebraic number field $K$ generated by $a$ and $e^a$, and we denote by $h$ its degree over the rational number field. If $\xi$ is any number in $K$, the norm $\mathfrak{N}(\xi) = \xi^{(1)} \dots \xi^{(h)}$ means the product of all $h$ conjugates $\xi^{(1)}, \dots, \xi^{(h)}$ of $\xi$, and we define
\[ \overline{|\xi|} = \max (|\xi^{(1)}|, \ldots, |\xi^{(h)}|). \]

Take $m=h+1$, $\rho_k=(k-1)a$ ($k=1,\dots,m$), $n_1=n_2=\dots=n_m=n \geq 1$. The number $n$ will be arbitrarily large, and we shall denote by $c_1, c_2, \dots$ positive rational integers which are independent of $n$.

Now the number $P_k(1)$ lies in $K$. Using\eqref{eq:27} for the $h$ conjugate fields we obtain from\eqref{eq:32} that
\begin{equation} \overline{|P_{k}(1)|} < c_{1}^{n}, \tag{33}\label{eq:33} \end{equation}
since $N=m(n+1)-1=mn+h < c_2 n$. Also
\[ R(1) = P_1(1) e^{\rho_1} + \dots + P_m(1) e^{\rho_m} \]
lies in $K$. In virtue of our former estimate for the denominator of $P_k(1)$ we can determine a number
\begin{equation} T = c_3^n n! \tag{34}\label{eq:34} \end{equation}
such that
\[ \xi = TR(1) \]
is an \textit{integer} in $K$. By\eqref{eq:30}, this integer is not 0; therefore
\begin{equation} |\mathfrak{N}(\xi)| \geq 1. \tag{35}\label{eq:35} \end{equation}

On the other hand, by\eqref{eq:29} and\eqref{eq:34}, we have
\[ |\xi| < c_4^n(n!)^{-h} \]
for \textit{one} conjugate, $\xi = \xi^{(1)}$ say, but not necessarily for any other conjugate $\xi^{(2)}, \ldots, \xi^{(h)}$, since the $h-1$ numbers $e^x$ ($x=\rho_k^{(2)},\ldots,\rho_k^{(h)}$) are perhaps not conjugate to $e^{\rho_k}$. However, by\eqref{eq:33} and\eqref{eq:34}, we obtain for \textit{all} conjugates of $\xi$ the estimate
\[ \overline{|\xi|} < c_5^n n!. \]
Hence
\[ |\mathfrak{N}(\xi)| < c_4^n(n!)^{-h}(c_5^n n!)^{h-1} = \frac{c_6^n}{n!} < 1, \]
for all sufficiently large $n$, contradictory to\eqref{eq:35}.

This proves that $e^a$ is a transcendental number for all real algebraic $a \neq 0$. In particular, the number $e$ itself is transcendental.

\section{The determinant of m approximation forms}

The essential point in the preceding proof was the fact that the algebraic integer $\xi$ is not 0, and this followed from\eqref{eq:30}. If $\rho_1, \dots, \rho_m$ again are arbitrary different complex numbers, we can no longer assert that $R(1) \neq 0$. To overcome this difficulty we proceed in the following way.

For any fixed $k=1, 2, \dots, m$ we choose $n_l = n \geq 1$ ($l=1, \dots, k$) and $n_l = n-1$ ($l=k+1, \dots, m$) and denote the corresponding approximation form by
\[ R_{k} = P_{k1} e^{\rho_{1}x} + \dots + P_{km} e^{\rho_{m}x} \quad (k=1,\dots,m). \]
The degree of the polynomial $P_{kl}$ is $n_l$, and this number equals $n$ or $n-1$, according to $k \geq l$ or $k < l$. Now consider the determinant $\Delta = \Delta(x)$ of the $P_{kl}$. Its $m!$ terms are polynomials in $x$, and all of them have a degree $< mn$, except the term corresponding to the main diagonal which has the degree $mn$. It follows that the polynomial $\Delta$ has the exact degree $mn$.

Denoting the minors of the elements of the first column of $\Delta$ by $\Delta_1, \dots, \Delta_m$ we have
\begin{equation} \Delta = (\Delta_1 R_1 + \dots + \Delta_m R_m) e^{-\rho_1 x}. \tag{36}\label{eq:36} \end{equation}
The function $R_k$ vanishes at $x=0$ of order
\[ (n_1 + \dots + n_m) + m-1 = mn + k-1 \geq mn; \]
therefore $\Delta$ vanishes at $x=0$ at least of order $mn$. This proves that
\[ \Delta(x) = \gamma x^{mn}, \quad \Delta(1) = \gamma \neq 0. \]

Now\eqref{eq:36} implies that at least one of the $m$ numbers $R_k(1)$ ($k=1,\dots,m$) is not 0, and this suffices, in particular, for the extension of the proof of \S 10 to the case of a complex algebraic $a \neq 0$.

\section{Algebraic independence}

Suppose that $p$ algebraic numbers $a_1, \dots, a_p$ are related by a homogeneous linear equation $g_1 a_1 + \dots + g_p a_p = 0$ with rational integral coefficients $g_1, \dots, g_p$, not all 0. Then the $p$ numbers $\eta_k = e^{a_k}$ ($k=1,\dots,p$) satisfy the algebraic equation $\eta_1^{g_1} \dots \eta_p^{g_p} = 1$ with rational coefficients. We are going to prove the converse: If $a_1, \dots, a_p$ are algebraic numbers such that $g_1 a_1 + \dots + g_p a_p \neq 0$ for all rational integers $g_1, \dots, g_p$ except for $g_1=0, \dots, g_p=0$, then the $p$ numbers $e^{a_1}, \dots, e^{a_p}$ are not related by an algebraic equation with algebraic coefficients. This means in case $p=1$ that $e^{a_1}$ is transcendental for algebraic $a_1 \neq 0$. In particular, $e$ is transcendental, since $1 \neq 0$.

Assume that the statement is not true; then there exists a polynomial $G = G(y_1, \dots, y_p)$ in $p$ variables $y_1, \dots, y_p$ and with integral algebraic coefficients, not all 0, such that $G$ vanishes for $y_k = e^{a_k} \quad (k=1,\dots,p)$. Let $d$ be the total degree of $G(y_1, \dots, y_p)$ and denote by $h$ the degree of the algebraic number field $K$ generated by $a_1, \dots, a_p$ and the coefficients of $G$. We choose a rational integer $f > d$ such that
\begin{equation} \prod_{k=1}^{p} (f-d+k) > \left(1-\frac{1}{h}\right) \prod_{k=1}^{p} (f+k); \tag{37}\label{eq:37} \end{equation}
this is possible because the difference of these two polynomials in $f$ has a positive highest coefficient, namely $\frac{1}{h}$.

The number of all monomials $y_1^{g_1} \dots y_p^{g_p}$ of total degree $\leq f$ or $\leq f-d$ is $m = \binom{f+p}{p}$ or $r = \binom{f-d+p}{p}$, respectively. Denoting them by $Y_1, \dots, Y_m$ and $Z_{m-r+1}, \dots, Z_m$, we obtain
\[ Z_k G = \alpha_{k1} Y_1 + \dots + \alpha_{km} Y_m \quad (k=m-r+1,\dots,m), \]
where $\alpha_{kl}$ is 0 or a coefficient of $G$. It is clear that the $r$-rowed coefficient matrix $(\alpha_{kl})$ has the rank $r$, since the polynomials $Z_k G$ are not related by any homogeneous linear equation with constant coefficients, not all 0.

Let
\[ Y_l = y_1^{g_{l1}} \dots y_p^{g_{lp}}, \quad \rho_l = g_{l1} a_1 + \dots + g_{lp} a_p \quad (l=1,\dots,m). \]
Then $\rho_1, \dots, \rho_m$ are $m$ different numbers in $K$ and
\[ \alpha_{k1} e^{\rho_1} + \dots + \alpha_{km} e^{\rho_m} = 0 \quad (k=m-r+1,\dots,m). \]

Now consider the $m$ approximation forms
\[ R_k(x) = P_{k1}(x) e^{\rho_1 x} + \dots + P_{km}(x) e^{\rho_m x} \quad (k=1,\dots,m) \]
of \S 11. We know that the $m$-rowed matrix $(P_{kl}(1))$ has a non-vanishing determinant $\Delta(1)$, so we can find $m-r$ rows, say for $k=k_1, \dots, k_{m-r}$, which together with the $r$ rows of the matrix $(\alpha_{kl})$ constitute a non-vanishing determinant. Write
\[ P_{k_t l}(1) = \alpha_{tl}, \quad R_{k_t}(1) = \beta_t \quad (t=1,\dots,m-r); \]
then
\[ \alpha_{k1} e^{\rho_1} + \dots + \alpha_{km} e^{\rho_m} = \beta_k \quad (k=1,\dots,m) \]
where $\beta_k = 0$ for $k=m-r+1, \dots, m$. Let $\Delta$ be the determinant of the $\alpha_{kl}$ and $\Delta_1, \dots, \Delta_m$ the minors corresponding to the first column. It follows that
\begin{equation} 0 \neq \Delta = (\Delta_1 \beta_1 + \dots + \Delta_{m-r} \beta_{m-r}) e^{-\rho_1}. \tag{38}\label{eq:38} \end{equation}

Now we shall apply the results of \S 9. Since $n_l = n$ or $n-1$ we can determine a number
\[ T = c_7^n n! \]
such that the $(m-r)m$ numbers $T\alpha_{kl}$ ($k=1,\dots,m-r; l=1,\dots,m$) are integers in $K$. This implies that also $T^{m-r}\Delta$ is integral, whence
\begin{equation} |\mathfrak{N}(\Delta)| \geq T^{h(r-m)} = c_8^{-n}(n!)^{h(r-m)}. \tag{39}\label{eq:39} \end{equation}

Moreover, by\eqref{eq:32},
\begin{equation} \overline{|\alpha_{kl}|} < c_9^n \quad (k=1,\dots,m; l=1,\dots,m), \tag{40}\label{eq:40} \end{equation}
\begin{equation} \overline{|\Delta|} < c_{10}^n. \tag{41}\label{eq:41} \end{equation}

On the other hand, by\eqref{eq:29},
\[ |\beta_t| < c_{11}^n (n!)^{-m} < c_{12}^n (n!)^{-m} \quad (t=1,\dots,m-r), \]
so that, by\eqref{eq:38} and\eqref{eq:40},
\begin{equation} |\Delta| < c_{13}^n (n!)^{-m}. \tag{42}\label{eq:42} \end{equation}

The estimates\eqref{eq:41} and\eqref{eq:42} imply
\[ |\mathfrak{N}(\Delta)| < c_{14}^n (n!)^{-m}. \]

Letting $n \to \infty$ and comparing with\eqref{eq:39} we see that
\[ m \leq h(m-r), \]
\[ r \leq \left(1-\frac{1}{h}\right)m. \]
Because of the definitions of $m$ and $r$, this contradicts\eqref{eq:37}.

The result can be formulated in another way. Suppose that $b_1, \dots, b_r$ are different algebraic numbers. If $p$ of these $r$ numbers, and not more, are linearly independent in the field of rational numbers, then we can find $p$ linearly independent algebraic numbers $a_1, \dots, a_p$, such that
\[ b_k = g_{k1}a_1 + \dots + g_{kp}a_p \quad (k=1,\dots,r) \]
with rational integral coefficients $g_{kl}$. Now consider the rational function
\[ f(y_1, \dots, y_p) = \sum_{k=1}^r c_k y_1^{g_{k1}} \dots y_p^{g_{kp}} \]
of the variables $y_1, \dots, y_p$ with arbitrary algebraic coefficients $c_1, \dots, c_r$, not all 0. It cannot vanish identically in $y_1, \dots, y_p$, since the sequences of exponents $g_{k1}, \dots, g_{kp}$ are different, for $k=1, \dots, r$. Now our result shows that $f$ cannot vanish for $y_1=e^{a_1}, \dots, y_p=e^{a_p}$; hence
\begin{equation} c_1 e^{b_1} + \dots + c_r e^{b_r} \neq 0. \tag{43}\label{eq:43} \end{equation}

This is the Lindemann-Weierstrass theorem: If $b_1, \dots, b_r$ are different algebraic numbers, then $e^{b_1}, \dots, e^{b_r}$ are not related by a homogeneous linear equation with algebraic coefficients, not all 0.

Vice versa, this theorem contains our previous result, since an algebraic equation between $e^{a_1}, \dots, e^{a_p}$ with algebraic coefficients would mean that a certain polynomial $f(y_1, \dots, y_p)$ vanishes for $y_1 = e^{a_1}, \dots, y_p = e^{a_p}$. Since $a_1, \dots, a_p$ are linearly independent in the rational number field, we would obtain a contradiction to\eqref{eq:43}.

\section{Another expression for the remainder term $R(x)$}

We shall devote the rest of this chapter to study more closely the analytical properties of approximation forms. We start from the complex integral
\begin{equation} J = \frac{1}{2\pi i} \int_C \frac{e^{xz}}{Q(z)} dz, \quad Q(z) = \prod_{k=1}^{m} (z - \rho_k)^{n_k+1}, \tag{44}\label{eq:44} \end{equation}
where the $\rho_k$ and $n_k$ have their former meaning and $C$ is a simply closed curve in positive direction which contains $\rho_1, \ldots, \rho_m$ in its interior. Inserting
\[ e^{xz} = e^{x\rho_k} \sum_{l=0}^{\infty} \frac{x^l (z-\rho_k)^l}{l!} \]
we see that the residue of the integrand at $z=\rho_k$ takes the form $Q_k e^{\rho_k x}$, where $Q_k=Q_k(x)$ is a polynomial in $x$ of degree $\leq n_k$. Then, by the residue theorem,
\[ J = Q_1 e^{\rho_1 x} + \dots + Q_m e^{\rho_m x}. \]

On the other hand,
\[ J = \sum_{l=0}^{\infty} a_l \frac{x^l}{l!}, \quad a_l = \frac{1}{2\pi i} \int_C \frac{z^l}{Q(z)} dz. \]

If we take for $C$ a curve which contains the whole circles $|z| \leq |\rho_k|$ ($k=1,\dots,m$) in its interior, we may use the descending power series expansion
\[ \frac{1}{Q(z)} = z^{-N-1} \prod_{k=1}^{m} \left(1 - \frac{\rho_k}{z}\right)^{-n_k-1} = z^{-N-1} + \dots \]
This proves that
\[ a_l = 0 \quad (l=0,1,\dots,N-1), \quad a_N = 1. \]

Now the uniqueness theorem of \S 7 shows that $Q_k(x) = P_k(x)$ and $J = R(x)$.

The expression
\begin{equation} R(x) = \frac{1}{2\pi i} \int_C \frac{e^{xz}}{Q(z)} dz \tag{45}\label{eq:45} \end{equation}
of the remainder term as a simple complex integral is more elegant than the expression as an $(m-1)$-fold real integral in\eqref{eq:28}, but it is inconvenient if one wants to prove the results of \S 8. Without using the uniqueness theorem we can transform\eqref{eq:45} into\eqref{eq:28} in the following way. Suppose $x>0$, then $C$ can be deformed into a straight line $L$ from $c-i\infty$ to $c+i\infty$, where $c$ is a real number greater than the real parts of the $\rho_k$. Substitute
\[ (z-\rho_k)^{-n_k-1} = \frac{1}{n_k!} \int_0^\infty t_k^{n_k} e^{(\rho_k-z)t_k} dt_k \quad (k=1,\dots,m-1) \]
and interchange the order of integration; then
\[ R(x) = \int_0^{\infty} \dots \int_0^{\infty} \left\{ \frac{1}{2\pi i} \int_L \frac{e^{(x-t_1-\dots-t_{m-1})z}}{(z-\rho_m)^{n_m+1}} dz \right\} \prod_{k=1}^{m-1} \frac{t_k^{n_k}}{n_k!} e^{\rho_k t_k} dt_k. \]
But
\[ \frac{1}{2\pi i} \int_L \frac{e^{tz}}{(z-\rho_m)^{n_m+1}} dz = \begin{cases} \frac{t^{n_m}}{n_m!} e^{\rho_m t} & (t>0) \\ 0 & (t<0), \end{cases} \]
and\eqref{eq:28} follows immediately.

It is simple to get rid of the restriction $x>0$ in\eqref{eq:28}: Substitute $t_k x$ for $t_k$, then
\begin{equation} R(x) = x^N \int \dots \int_{\substack{t_1 + \dots + t_m = 1 \\ t_1 > 0, \dots, t_m > 0}} \prod_{k=1}^{m} \left( \frac{t_k^{n_k}}{n_k!} e^{\rho_k t_k x} \right) dt_1 \dots dt_{m-1}, \tag{46}\label{eq:46} \end{equation}
and this holds for arbitrary complex $x$. Since $R(x) = \frac{x^N}{N!} + \dots$, the integral in\eqref{eq:46} has for $x=0$ the value $\frac{1}{N!}$; hence
\[ |R(x)| \le \frac{|x|^N}{N!} e^{\rho|x|}, \qquad \rho = \max(|\rho_1|, \dots, |\rho_m|), \]
which is a refinement of\eqref{eq:29}.

It remains to prove formula\eqref{eq:27} for $P_k(x)$ from our present point of view. We know that $P_k(x)$ is the coefficient of $(z-\rho_k)^{n_k}$ in the power series expansion of the function
\[ f_k(z) = e^{x(z-\rho_k)} \prod_{\substack{l=1\\l\neq k}}^m (z-\rho_l)^{-n_l-1} \]
at the point $z=\rho_k$,
\[ P_k(x) = \frac{1}{n_k!} \{D_z^{n_k} f_k(z)\}_{z=\rho_k}. \]
Putting
\[ \prod_{\substack{l=1\\ l \neq k}}^{m} (z - \rho_l)^{-n_l-1} = g(z), \]
we obtain
\[ D_z^{n_k} f_k(z) = e^{x(z-\rho_k)} (x+D_z)^{n_k} g(z). \]

Now consider more generally the expression $\varphi(x+D_z)g(z)$, for any polynomial $\varphi(x)$. Then, by Taylor's formula,
\[ \varphi(x+D_z)g(z) = \sum_{l=0}^{\infty} \frac{D_x^l \varphi(x) D_z^l g(z)}{l!} = g(z+D_x)\varphi(x). \]
It follows that
\[ P_k(x) = g(\rho_k+D) \frac{x^{n_k}}{n_k!} = \prod_{\substack{l=1\\l\neq k}}^{m} (\rho_k-\rho_l+D)^{-n_l-1} \frac{x^{n_k}}{n_k!}; \]
q.e.d.

\section{The interpolation formula}

The integral $J$ of \S 13 also appears in the solution of the following interpolation problem: Let an analytic function $f(z)$ be regular in a domain $D$ of the complex $z$-plane, and let $n$ points $z_1, \ldots, z_n$ of $D$ be given; to determine a polynomial $H_{n-1}$ of degree $\leq n-1$ such that the fraction
\[ T_n(z) = \frac{f(z) - H_{n-1}(z)}{\prod_{k=1}^n (z - z_k)} \]
is regular in $D$.

It is clear that there cannot be more than one solution, since the difference of two solutions $H_{n-1}$ would be a polynomial of degree $< n$ and divisible by the polynomial $\prod_{k=1}^n (z - z_k)$ of degree $n$.

Put
\[ F_k(z) = (z-z_1) \dots (z-z_k) \quad (k=0,\dots,n), \]
so that
\begin{equation} \begin{aligned} F_k(z) &= (z-z_k)F_{k-1}(z) \quad (k=1,\dots,n), \\ (z-z_k)F_{k-1}(\zeta) - F_k(\zeta) &= (z-\zeta)F_{k-1}(\zeta), \\ \frac{1}{z-\zeta} \left( \frac{F_{k-1}(\zeta)}{F_{k-1}(z)} - \frac{F_k(\zeta)}{F_k(z)} \right) &= \frac{F_{k-1}(\zeta)}{F_k(z)}. \end{aligned} \tag{47}\label{eq:47} \end{equation}

Suppose that also $\zeta$ lies in $D$ and consider a simply closed curve $C$ in $D$ whose interior lies in $D$ and contains the $n+1$ points $z_1, \ldots, z_n, \zeta$. Define
\begin{equation} \begin{aligned} a_{k-1} &= \frac{1}{2\pi i} \int_C \frac{f(z)}{F_k(z)} dz, \\ G_k(\zeta) &= \frac{1}{2\pi i} \int_C \frac{F_k(\zeta)}{F_k(z)} \frac{f(z)}{z-\zeta} dz; \end{aligned} \tag{48}\label{eq:48} \end{equation}
then
\[ G_0(\zeta) = f(\zeta) \]
and, by\eqref{eq:47},
\[ G_{k-1}(\zeta) - G_k(\zeta) = a_{k-1}F_{k-1}(\zeta) \quad (k=1,\dots,n); \]
whence
\[ f(\zeta) = a_0 F_0(\zeta) + a_1 F_1(\zeta) + \dots + a_{n-1} F_{n-1}(\zeta) + G_n(\zeta). \]

Therefore the solution of the interpolation problem is given by
\[ \begin{aligned} H_{n-1}(\zeta) &= a_0 F_0(\zeta) + a_1 F_1(\zeta) + \dots + a_{n-1} F_{n-1}(\zeta), \\ T_n(\zeta) &= \frac{G_n(\zeta)}{F_n(\zeta)} = \frac{1}{2\pi i} \int_C \frac{f(z)}{F_n(z)} \frac{dz}{z - \zeta}. \end{aligned} \]

Now consider an infinite sequence of points $z_1, z_2, \dots$ in $D$ and suppose that $\lim_{n \to \infty} G_n(z) = 0$ for all $z$ in a domain $D_0 \subset D$. Then
\[ f(z) = a_0 F_0(z) + a_1 F_1(z) + \dots \quad (z \text{ in } D_0) \]
and it follows that $a_n \neq 0$ for infinitely many $n$ except when $f(z)$ is a polynomial.

We apply this to the special case $f(z) = e^{xz}$ and take for $z_1, \dots, z_{N+1}$ the set consisting of $n_k+1$ times the point $\rho_k$ ($k=1,\dots,m$). By\eqref{eq:44} and\eqref{eq:48}, we see that the coefficient $a_N$ is exactly the integral $J = R(x)$ of \S 13. In particular, choose $n_k = n$ ($k=1,\dots,r+1$) and $n_k = n-1$ ($k=r+2,\dots,m$) and take $n = 0, 1, \dots$; $r = 0, 1, \dots, m-1$, so that $N = mn+r = 0, 1, \dots$. It is easy to see from the expression for $G_N(z)$ that $\lim_{N \to \infty} G_N(z) = 0$ for all $z$. Since $e^{xz}$ is not a polynomial, this implies that $R(1) \neq 0$ for infinitely many $N$. The finer algebraical approach of \S 11 showed that for every $n \geq 0$ the interval $mn \leq N < m(n+1)$ contains at least one $N$ with $R(1) \neq 0$; but this fact is not needed for the proof of the transcendency of $e^a$ for algebraic $a \neq 0$, though it is important for the more general problem solved in \S 12.

\section{Concluding remarks}

It should be mentioned that the preceding proofs of the transcendency of $e$ and $\pi$ and of the algebraic independence of $e^{\alpha_1}, \dots, e^{\alpha_p}$, for linearly independent algebraic $\alpha_1, \dots, \alpha_p$, are not the simplest to be found in literature. Our proofs are related to the original work of Hermite; however, our procedure in constructing the approximation forms is somewhat more algebraic, and this has been decisive for the generalization which we shall investigate in the next chapter.

Two characteristic properties of the exponential function $y = e^x$ were used in our proofs, namely the differential equation $y' = y$ and the addition theorem $e^{x+t} = e^x e^t$. Our generalizations will go in two different directions: Either we consider solutions of linear differential equations without assuming an addition theorem, or we deal with functions satisfying an algebraic addition theorem. In the first case we are led to the problems considered in the second chapter. The second case brings us to the study of elliptic functions from the arithmetic point of view, in the last chapter, and of the function $a^x$ for algebraic $a \neq 0$, in the third chapter.
